% Problem record op_330054883db95e6b. This TeX file is derived from
% database/problems_json/op_330054883db95e6b.json by scripts/sync-tex.mjs; edit the JSON record
% and rerun the script rather than editing this file.
\section{Sublinear logical quantum space for RSA factoring}
\label{sec:op_330054883db95e6b}

\paragraph{Problem.}
Can every balanced RSA semiprime be factored in polynomial time using $o(n)$ simultaneously live logical qubits?

Let $N=pq$ be an $n$-bit integer whose distinct odd prime factors each have bit length $n/2+O(1)$. Does there exist a uniform hybrid algorithm satisfying

\begin{equation}
q(n)=o(n),\qquad T_{\mathrm{total}}(n)=n^{O(1)},
\label{eq:3300-1}
\end{equation}

and factoring every input with success probability at least $2/3$?

Here $q(n)$ counts all simultaneously live logical qubits, including coherent control registers, data, work qubits, quantum memory, and supplied resource states. Polynomial classical storage is permitted and the input $N$ may be classical. Intermediate measurements and resets are allowed. Runtime includes repetitions and classical processing.

A particularly concrete stronger milestone is $q(n)=O(n^{1-\delta})$ for some fixed $\delta>0$. The basic question only asks for $o(n)$.

The displayed definitions, constraints, and target bounds are recorded in Eqs.~\eqref{eq:3300-1}.

\subsection*{Status}

Unsolved

\subsection*{Source}

This precise formulation is editor wording based on the unresolved direction and limitations documented in the cited primary literature \sourcecite{ref:3300-chevignard25}{Chevignard25}\sourcecite{ref:3300-ragavan26}{Ragavan26}\sourcecite{ref:3300-kahanamokumeyer25}{KahanamokuMeyer25}; it is not presented as a verbatim conjecture of those authors.

\subsection*{Progress}

\begin{itemize}
  \item CRYPTO 2025, full version revised December 31, 2025. Chevignard, Fouque, and Schrottenloher factor $n$-bit RSA moduli with $n/2+o(n)$ logical qubits. Their modular-exponentiation technique uses sublinear work space, but the full algorithm still has linear total quantum width. These are different claims. \sourcecite{ref:3300-chevignard25}{Chevignard25}

  \item February 4, 2026. Ragavan and Vaikuntanathan give a space-efficient Regev implementation with $O(n\log n)$ qubits and $O(n^{3/2}\log n)$ gates per circuit. This improves the space–work balance of that approach but remains superlinear in the qubit parameter relevant here. \sourcecite{ref:3300-ragavan26}{Ragavan26}

  \item July 2, 2026 revision. For $P^2Q$ with $\log Q=\Theta(n^a)$ and $2/3<a<1$, the Jacobi construction has $\widetilde O(\log Q)=o(n)$ quantum space and depth. This is genuine progress for a classically difficult factoring family, but not for balanced RSA semiprimes. \sourcecite{ref:3300-kahanamokumeyer25}{KahanamokuMeyer25}

  \item Earlier “sublinear-resource” claims require care. A primary critique by Khattar and Yosri found failures in the proposed Schnorr/QAOA factoring approach even when idealizing its optimizer. Such heuristic proposals do not supply the uniform polynomial-time guarantee in this question. \sourcecite{ref:3300-khattar23}{Khattar23}

  \item Retained as open for RSA. A broader claim that no classically difficult factoring family admits proven sublinear quantum space would already be false.
\end{itemize}

\subsection*{References}

\begin{enumerate}[label={},leftmargin=4.8em,labelsep=0.5em]
  \footnotesize
  \item[\textup{[Chevignard25]}]\label{ref:3300-chevignard25}
  Clémence Chevignard, Pierre-Alain Fouque, and André Schrottenloher. \emph{Reducing the Number of Qubits in Quantum Factoring}. CRYPTO 2025; \href{https://eprint.iacr.org/2024/222}{ePrint 2024/222}, full version last revised December 31, 2025.

  \item[\textup{[Ragavan26]}]\label{ref:3300-ragavan26}
  Seyoon Ragavan and Vinod Vaikuntanathan. \emph{Space-Efficient and Noise-Robust Quantum Factoring}. \href{https://link.springer.com/article/10.1007/s00145-026-09572-x}{Journal of Cryptology 39, article 14 (2026)}, published February 4, 2026.

  \item[\textup{[KahanamokuMeyer25]}]\label{ref:3300-kahanamokumeyer25}
  Gregory D. Kahanamoku-Meyer, Seyoon Ragavan, Vinod Vaikuntanathan, and Katherine Van Kirk. \emph{The Jacobi Factoring Circuit: Quantum Factoring with Near-Linear Gates and Sublinear Space and Depth}. STOC 2025; \href{https://arxiv.org/abs/2412.12558v4}{arXiv:2412.12558v4}, July 2, 2026. Especially the input-family restriction and the sublinear-space result.

  \item[\textup{[Khattar23]}]\label{ref:3300-khattar23}
  Tanuj Khattar and Noureldin Yosri. \emph{A comment on “Factoring integers with sublinear resources on a superconducting quantum processor”}. \href{https://arxiv.org/abs/2307.09651v2}{arXiv:2307.09651v2} (2023), preprint.
\end{enumerate}

\subsection*{Comment}

This problem is a precise resource target formulated here from the quantum-space frontier.

The important distinction is between reducing workspace around an existing coherent register and reducing total coherent information storage. Replacing a large work register by recomputation does not resolve the problem when another register remains of size $\Theta(n)$.

There is no inference here that $\Omega(n)$ logical qubits are necessary for all polynomial-time hybrid factoring algorithms. Establishing such a universal lower bound would be a different, substantially stronger undertaking. The permitted classical memory also means that an information-counting argument based only on the $n$-bit input length is insufficient: the input need not reside in quantum memory.

\subsection*{Field}

Quantum algorithm

\subsection*{Topic}

Computational complexity and computability; Quantum circuit complexity

\subsection*{ID}

\texttt{op\_330054883db95e6b}
